\section{观察、干预与反事实回答不同问题}
\label{sec:13-Machine-Learning/12-Causal-Inference/01}

数据告诉你：接受新药的患者康复率比没接受的高 15\%。能直接说药有效吗？

不能。也许病情更重的人更倾向于接受治疗——而病情本身就是康复率的最强预测因素。你看到的 15\% 是治疗效果和治疗前健康状况的混合：治疗确实有帮助，但接受治疗的人起点更低，观察到的差距被压缩了。极端情况下，即使药效显著，治疗组的观察结果也可能更差（如果选择偏差足够强）。

也可能是反方向：康复希望大的患者，医生才敢用新药。于是治疗组看起来更好，不是药有用，而是选进去的人本来就更可能康复。

从同一批数据中"看见某些人接受了治疗"、"主动让某些人接受治疗"、追问"同一个人如果没接受治疗会怎样"，是三个完全不同的问题。变量名字可以完全一样——\(X\) 表示治疗，\(Y\) 表示结果——但数据是怎样产生的不同。不先说清你究竟想比较什么，回归系数就不会自动翻译成因果效应。

\subsection{三个问题共享变量，却不共享答案}

第一个问题：已经观察到某人接受了治疗，他的结果怎样分布？条件概率回答：

\[
P(Y\mid X=x).
\]

这是统计学最基础的工具。它描述的是在被动的观察世界里，看到 \(X=x\) 的同时 \(Y\) 如何分布。"筛出恰好满足某条件的人，看看他们的结果。"

第二个问题：如果外部强制每个人都接受治疗——不管他们自己愿不愿意、不管医生怎么判断——结果会怎样分布？记作：

\[
P\bigl(Y\mid \operatorname{do}(X=x)\bigr).
\]

\(\operatorname{do}(\cdot)\) 是 Pearl 引入的记号。它表示切断了 \(X\) 原有的生成机制，从外部把它设为 \(x\)。不是"等待自然发生然后筛选"，而是"动手改变然后观察"。

第三个问题：这个具体的人实际接受了治疗并且康复了。如果当初没给他治疗，他会康复吗？记作：

\[
P\bigl(Y(0)=1\mid X=1,Y=1\bigr).
\]

\(Y(0)\) 是同一单位在未接受治疗时的潜在结果。这个问题既用了事实信息（他真的接受了治疗、真的康复了），又追问了未发生的世界（如果没治疗呢）。它通常比总体干预分布需要更强的模型——不仅要推断治疗在总体中的平均效果，还要推断特定个体的"如果"。

三个问题用同一个数据集可能给出三个完全不同的答案。不是计算出了错，是问题本身就不同。

\subsection{条件化是在筛选，干预是在改写机制}

观察 \(X=x\) 做条件化，是从原始总体中筛出一个子群：所有自然满足 \(X=x\) 的人。这个操作不改变任何人的 \(X\) 值，不改变任何人的 \(Y\) 值，只是选择性地看一部分人。

\(\operatorname{do}(X=x)\) 是改写生成机制：把决定 \(X\) 的那个方程从"由原因和噪声产生"替换为"恒等于 \(x\)"。其余机制——从 \(X\) 到 \(Y\) 的那条因果路径、其他变量的生成方式——保持不变。

什么时候两种操作给出相同的分布？当 \(X\) 的赋值近似随机（随机对照试验），或所有同时影响 \(X\) 和 \(Y\) 的变量都已被适当控制住了（后门准则满足）。如果不满足——两个分布可以差得天远。

一个经典的例子：气压计读数与暴风雨。气压低，暴风雨概率高——条件概率说 \(P(\text{暴风雨}\mid\text{低读数})\) 很大。但如果你强行把气压计指针拨到"低"，天气会变吗？不会。

\[
P(Y\mid X=x)
\ne
P\bigl(Y\mid \operatorname{do}(X=x)\bigr).
\]

这个不等式不是数学上的细枝末节，是整个分析的岔路口。数据告诉你：看到低读数之后下雨更频繁。但你绝对不能拿这个数字来指导"要不要拨指针"。两个数字对应两个不同的行动——带伞 vs 拨指针——混淆它们会做出荒唐的决策。记号 \(\operatorname{do}\) 正是把这两种操作显式分开。

\subsection{先定义要比较的总体量}

对二元处理 \(X\in\{0,1\}\)，总体平均处理效应：

\[
\tau
=
\E\bigl[Y(1)-Y(0)\bigr].
\]

若干预分布可识别，也等于：

\[
\tau
=
\E\bigl[Y\mid\operatorname{do}(X=1)\bigr]
-
\E\bigl[Y\mid\operatorname{do}(X=0)\bigr].
\]

这是最常见的报告数字，但远不是唯一值得关心的。已处理人群的平均效应：

\[
\E\bigl[Y(1)-Y(0)\mid X=1\bigr]
\]

回答的是"实际接受了治疗的那些人，他们受益了多少"。这个数字可能和 \(\tau\) 完全不同——如果治疗分配跟预期获益相关的话。还有特定特征 \(C=c\) 下的条件效应、某项政策在目标人群中的总体价值。它们的目标人群不同、加权方式不同，即使来自同一个研究也不该混在一起报一个数字。

定义了 estimand，才能谈识别——需要什么假设、什么数据。不先定义 estimand，直接跑回归只会得到一个系数；那个系数回答的是什么问题，回归不会替你判断。

\subsection{预测准确不等于干预正确}

回到急诊分诊的场景。令 \(W\) 表示患者特征，\(A\in\{0,1\}\) 表示是否加强监护，\(Y=1\) 表示二十四小时内严重恶化。预测模型估计：

\[
r(w)=P(Y=1\mid W=w),
\]

把患者按风险从高到低排序。这很有用——知道哪些患者最有可能恶化，可以把有限的监护资源往他们倾斜。

但如果加强监护的目的不是预测恶化，而是通过早期干预防止恶化，资源分配真正该看的是：

\[
b(w)=\E\bigl[Y(0)-Y(1)\mid W=w\bigr],
\]

也就是"监护能降低多少恶化风险"。\(r(w)\) 和 \(b(w)\) 的排序可以完全不同。病情极重的患者可能 \(r(w)\) 极高——不管监护不监护都大概率恶化——\(b(w)\) 却趋近于零。另一位中等病情的患者，恶化风险中等，但监护能显著降低它——\(b(w)\) 很大。风险最高的不一定是获益最大的。把资源按风险排序分配，可能把大量资源花在干预无效的人身上，而真正能从干预中受益的人被漏掉。

历史数据还会把这个混淆锁死。高风险患者更常被医生安排监护，观察到的 \(Y\) 已经是干预后的结果。如果治疗有效，他们后来的恶化率已经比不治疗时低了。用这些标签训练"未经干预时的风险预测模型"，会把治疗成功误读为患者本来安全——模型看着很准，学到的是"被治疗者恶化少"，而不是"恶化风险低"。

行动改变了结果生成机制。增加样本量不能自动恢复未发生的反事实——再多的数据也是在干预之后的世界上采集的。一个变量有没有用，不止取决于它和结果的相关性，还取决于它产生在干预之前还是之后、位于哪条因果路径上。治疗后的中介变量预测力极强，但如果目标是估计治疗总效应，控制它就是截断了因果路径，改变了估计量的含义。

反过来也一样：因果效应估计未必适合个体预测。平均处理效应接近零，可能是因为所有个体的效应都小，也可能是正负效应在不同人群中互相抵消——药对一半人救命、对另一半人致命。只看平均，会错过最关键的异质性。先写清楚要估计的量，才知道需要什么数据和假设。直接跑回归只会吐出一个系数，不会替你决定那个系数回答了什么。
